飓风后鸟类在不同环境中的出现及占域状态变化，是确定监测重点的重要依据。本研究以波多黎各绿蜂鸟（\Rm）为对象，利用 2013--2025 年的 \nChecklists{} 份 eBird 清单，结合海拔与植被绿度分层、动态占域模型及空间和时间留出预测，分析飓风相关时期的生境利用差异，并比较不同调查记录筛选方案的环境覆盖。María 前、两次飓风之间及 Fiona 后的总体报告率分别为 7.59\%、4.62\% 和 6.21\%。标准化预测中，高海拔高绿度环境保持较高报告水平，后期绝对增幅的点估计也较大；低地部分环境的报告水平仍低于初期，但各环境层末期相对初期差值的区间均跨零。校正后的动态占域模型显示，高海拔高绿度网格的平均占域持续概率为 96.8\%，低海拔低、高绿度网格分别为 24.2\% 和 71.3\%。该概率指原已占据网格在下一主要时期仍被占据的概率。绿度与占域持续性的关联比其与定殖的关联更明确，表明低地整体报告偏低可以与较绿网格中较高的占域持续概率并存。预测改善主要集中于低、中海拔稀少报告的识别。相同清单筛选预算下，海拔分层排序将低地正报告覆盖率从 34.6\% 提高至 45.6\%，总体覆盖率则从 46.7\% 降至 35.2\%。仅优先选择高报告得分记录，难以兼顾不同环境的监测需求。后续调查应保留山地固定复访点，在低地覆盖不同绿度条件，并结合层内排序补充候选地点。调查应分别记录后续时期是否再次报告该物种，并在考虑不完全探测后估计占域持续概率和定殖概率，以检验环境差异及其时期变化。
